%% kapa: exam=1 round=32 year=2021 period=1 form=A subject=economics subject_ko=경제학원론 seq=11 number=51
%% source: Q-Net 감정평가사 기출문제 (cst00309) artlSeq=5208397 file=extracted/1008/5208397/A형/2021년도 제32회 감정평가사 1차 1교시 A형.pdf
%% conversion: kordoc 4.24.51 -> tools/gen1.py (auto); review=proofread-sonnet; flags=none
\begin{question}{51}
표는 기업 甲과 乙의 초기 보수행렬이다. 제도 변화 후, 오염을 배출하는 乙은 배출 1톤에서 2톤으로 증가하는데 甲에게 보상금 5를 지불하게 되어 보수행렬이 변화했다. 보수행렬 변화 전, 후에 관한 설명으로 옳은 것은? (단, 1회성 게임이며, 보수행렬 (\ \ ) 안 왼쪽은 甲, 오른쪽은 乙의 것이다.)
\par{\small\centering
\begin{tabular}{|c|c|c|c|}
\cline{3-4}
\multicolumn{2}{c|}{} & \multicolumn{2}{c|}{乙} \\ \cline{3-4}
\multicolumn{2}{c|}{} & 1톤 배출 & 2톤 배출 \\ \hline
\multirow{2}{*}{甲} & 조업중단 & (0, 4) & (0, 8) \\ \cline{2-4}
 & 조업가동 & (10, 4) & (3, 8) \\ \hline
\end{tabular}\par}
\begin{choices}
  \item 초기 상태의 내쉬균형은 (조업중단, 2톤 배출)이다.
  \item 초기 상태의 甲과 乙의 우월전략은 없다.
  \item 제도 변화 후 甲의 우월전략은 있으나 乙의 우월전략은 없다.
  \item 제도 변화 후 甲과 乙의 전체 보수는 감소했다.
  \item 제도 변화 후 오염물질의 총배출량은 감소했다.
\end{choices}
\end{question}
\answer{5}
